% ===== PRÉAMBULE CANONIQUE — à coller VERBATIM en tête de CHAQUE .tex =====
\documentclass[11pt,a4paper]{article}
\usepackage[utf8]{inputenc}\usepackage[T1]{fontenc}\usepackage{lmodern}
\usepackage[french]{babel}
\usepackage{amsmath,amssymb,mathtools}\usepackage{icomma}
\usepackage{siunitx}\sisetup{locale=FR}  % \qty{}{}, \si{}, \ang{}
\usepackage{xcolor}\usepackage{colortbl}
\usepackage{tikz}\usepackage{pgfplots}\pgfplotsset{compat=1.18}
\usepackage[siunitx,europeanresistors]{circuitikz} % circuits (élec.)
\usepackage[most]{tcolorbox}\usepackage{enumitem}\usepackage{array,booktabs}
\usepackage[a4paper,margin=2.2cm]{geometry}
\usetikzlibrary{arrows.meta,calc,angles,quotes,positioning,patterns,%
  backgrounds,shapes.geometric,decorations.pathmorphing,decorations.markings,intersections}
\definecolor{primary}{RGB}{222,49,99}\definecolor{primarydark}{RGB}{150,22,55}
\definecolor{defcol}{RGB}{0,114,178}\definecolor{methcol}{RGB}{52,92,148}
\definecolor{excol}{RGB}{0,135,90}\definecolor{attncol}{RGB}{200,40,55}
\tcbset{boxbase/.style={enhanced,breakable,boxrule=0.9pt,arc=2.5pt,
  left=3mm,right=3mm,top=2.3mm,bottom=2.3mm,fonttitle=\bfseries,coltitle=white}}
% --- boîtes TUTORIEL (noms EXACTS, ne pas renommer) ---
\newtcolorbox{cle}[1][]{boxbase,colback=defcol!6,colframe=defcol,
  title={\textsc{À retenir}\ifx\relax#1\relax\relax\else\textmd{ : \itshape #1}\fi}}
\newtcolorbox{methode}[1][]{boxbase,colback=methcol!7,colframe=methcol,sharp corners,
  title={\textsc{Méthode}\ifx\relax#1\relax\relax\else\textmd{ --- \itshape #1}\fi}}
\newtcolorbox{exemplebox}[1][]{boxbase,colback=excol!5,colframe=excol,
  title={\textsc{Exemple}\ifx\relax#1\relax\relax\else\textmd{ : \itshape #1}\fi}}
\newtcolorbox{piege}[1][]{boxbase,colback=attncol!6,colframe=attncol,
  title={\textsc{Piège}\ifx\relax#1\relax\relax\else\textmd{ : \itshape #1}\fi}}
\newtcolorbox{remarque}[1][]{boxbase,colback=black!4,colframe=black!45,
  title={\textsc{Remarque}\ifx\relax#1\relax\relax\else\textmd{ : \itshape #1}\fi}}
% --- boîtes EXERCICE / CORRIGÉ (noms EXACTS, auto-comptées) ---
\newtcolorbox[auto counter]{exo}[1][]{boxbase,colback=primary!4,colframe=primary,
  title={\textsc{Exercice}~\thetcbcounter\ifx\relax#1\relax\relax\else\textmd{ --- \itshape #1}\fi}}
\newtcolorbox[auto counter]{corrige}[1][]{boxbase,colback=excol!4,colframe=excol,
  title={\textsc{Corrigé}~\thetcbcounter\ifx\relax#1\relax\relax\else\textmd{ --- \itshape #1}\fi}}
\newcommand{\vv}[1]{\vec{#1}}\newcommand{\ds}{\displaystyle}
\newcommand{\R}{\mathbb{R}}\newcommand{\N}{\mathbb{N}}\newcommand{\Z}{\mathbb{Z}}
% (un \usepackage{fancyhdr}+en-tête est possible ; il est ignoré côté web)
% =========================================================================
\begin{document}
% THEME_TITLE: Chute libre et tir horizontal
\sisetup{per-mode=symbol}

\noindent Un corps soumis à la \emph{seule} pesanteur tombe avec l'accélération $g$, la \textbf{même
pour tous les corps} quelle que soit leur masse. On étudie d'abord la chute verticale, puis le tir
horizontal, qui n'est que la superposition d'un MRU et d'une chute.

\begin{cle}[chute libre verticale (lâché sans vitesse)]
Avec $g\approx\qty{9,81}{\metre\per\second\squared}$ (souvent arrondi à \qty{10}{\metre\per\second\squared}) :
\[ h=\tfrac12 g\,t^2,\qquad v=g\,t,\qquad t=\sqrt{\frac{2h}{g}},\qquad v=\sqrt{2gh}. \]
La masse n'intervient pas : une plume et une bille tombent identiquement dans le vide.
\end{cle}

\begin{cle}[choix du signe et lancer vertical]
Le signe de $g$ dépend de l'orientation de l'axe : axe \emph{vers le bas} $\Rightarrow a=+g$ ;
axe \emph{vers le haut} $\Rightarrow a=-g$. Un corps lancé vers le haut à $v_0$ monte, s'arrête
($v=0$) à l'apogée, puis redescend ; sa hauteur maximale (depuis le point de lancer) vaut
$\;y_{\max}=\dfrac{v_0^2}{2g}$.
\end{cle}

\begin{cle}[tir horizontal : deux mouvements indépendants]
Un corps lancé horizontalement à $v_{0x}$ combine :
\[ \underbrace{x(t)=v_{0x}\,t}_{\text{MRU selon }x}\qquad\text{et}\qquad
   \underbrace{y(t)=\tfrac12 g\,t^2,\quad v_y=g\,t}_{\text{chute libre selon }y},\qquad
   v=\sqrt{v_x^2+v_y^2}. \]
Le \textbf{temps de vol ne dépend que de la hauteur} de chute, pas de $v_{0x}$.
\begin{center}
\begin{tikzpicture}[scale=0.95,>=Stealth]
  \draw[->,thick,black!70] (-0.2,0)--(5.4,0) node[right]{$x$};
  \draw[->,thick,black!70] (0,0.2)--(0,-3.6) node[below]{$y$};
  \node[above left,font=\scriptsize] at (0,0){$O$};
  \draw[primary,line width=1.1pt,domain=0:4.6,smooth,variable=\x] plot ({\x},{-0.17*\x*\x});
  \foreach \x in {0,1.6,2.8,3.8}{
     \pgfmathsetmacro\y{-0.17*\x*\x}
     \fill (\x,\y) circle (1.3pt);
     \draw[->,defcol,line width=0.8pt] (\x,\y)--++(0.9,0);
     \pgfmathsetmacro\vyy{-0.33*\x-0.15}
     \draw[->,attncol,line width=0.8pt] (\x,\y)--++(0,\vyy);}
  \node[defcol,font=\scriptsize] at (1.05,0.28){$v_x$ (cte)};
  \node[attncol,font=\scriptsize] at (4.4,-2.1){$v_y$ croît};
\end{tikzpicture}
\end{center}
\end{cle}

\begin{methode}[résoudre un tir horizontal]
\begin{enumerate}[label=\arabic*.,leftmargin=1.6em,topsep=1pt,itemsep=0pt]
  \item Temps de chute depuis la hauteur $h$ : $t=\sqrt{2h/g}$.
  \item Portée horizontale : $x=v_{0x}\times t$.
  \item Vitesse à l'impact : $v_x=v_{0x}$, $v_y=g\,t$, puis $v=\sqrt{v_x^2+v_y^2}$.
\end{enumerate}
\end{methode}

\begin{exemplebox}[lâché de \qty{125}{\metre}]
$h=\tfrac12 g t^2 \Rightarrow t=\sqrt{2\times125/10}=\qty{5}{\second}$ ; $v=g\,t=\qty{50}{\metre\per\second}$.
La vitesse gagne \qty{10}{\metre\per\second} chaque seconde : $0,10,20,30,40,\qty{50}{\metre\per\second}$.
\end{exemplebox}

\begin{piege}[le temps de chute est indépendant de la vitesse horizontale]
Une bille \emph{lâchée} et une bille \emph{lancée horizontalement} depuis la même hauteur touchent
le sol \textbf{en même temps} : $v_{0x}$ n'agit pas sur la chute. De même, $g$ ne dépend pas de la masse.
\end{piege}

\begin{remarque}[apogée au-dessus du sol]
Un objet lâché (ou lancé) d'une hauteur $h_0$ avec une vitesse ascensionnelle $v_0$ monte encore de
$v_0^2/2g$ : sa hauteur maximale par rapport au sol est $h_0+\dfrac{v_0^2}{2g}$.
\end{remarque}

\end{document}
